
\section{Introduction}
\label{sec:intro}

Machine translation (MT) into under-resourced target languages remains difficult even as neural and large-language-model (LLM) systems improve for high-resource pairs~\citep{koehn-knowles-2017-six,maillard-etal-2023-small}.
Interslavic (ISV; also Med\v{z}uslovjansky) is a constructed Slavic auxiliary language designed for mutual intelligibility across Slavic speech communities.
Relative to national Slavic languages with large web corpora and parallel datasets, ISV occupies a low-resource niche: usable dictionaries and community materials exist, but parallel Polish--ISV training data suitable for conventional neural MT are scarce, and automated evaluation infrastructure is not standardized in the same way as for WMT high-resource directions.

For Slavic speakers, Polish-to-ISV translation is practically interesting because Polish is a high-resource national language while ISV is intended as a bridge form.
For NLP methodology, the setting is interesting because it combines (i)~a target language with limited bitext, (ii)~available monolingual or curated target-side resources, and (iii)~commercial LLMs that can be prompted without gradient updates.
The methodological risk is equally clear: an attractive prompting trick can be mistaken for a validated translation method if experiments do not separate instruction-only behaviour from corpus-conditioned behaviour, and if they do not test whether apparent gains depend on thematic overlap between the source and the supplied corpus.

Prompting LLMs with instructions and in-context examples has become a practical route to translation without fine-tuning~\citep{brown2020language,hendy2023good,jiao2023chatgpt}.
A natural engineering idea for ISV is to place authentic ISV text in the prompt---here called \emph{corpus priming}---so that the model can condition generation on attested target-language material rather than on an instruction alone.
That idea is appealing, but it is not self-evidently beneficial.
Corpus material may encourage useful lexical and orthographic alignment; it may also encourage motif copying, register drift, or metric gains that do not correspond to communicative quality.
Moreover, if a source text already overlaps thematically with the priming corpus, apparent ``priming gains'' may partly reflect overlap rather than a generalizable priming intervention.

This paper reports EXP-004 Phase~2B: a controlled Polish-to-ISV study that separates Direct translation (no corpus) from Primed translation (authentic ISV corpus first, then translation), using identical model configurations, three independent repeats, and three source regimes that vary corpus overlap and genre:
\begin{itemize}[leftmargin=*]
  \item \textbf{HIGH}: narrative fiction deliberately inspired by the priming corpus;
  \item \textbf{LOW}: narrative fiction with intentionally weak thematic overlap;
  \item \textbf{UNSEEN}: shorter technical/expository educational text on human-voice biophysics, from a domain absent from the priming corpus registers.
\end{itemize}
The primary quantitative quantity is the paired difference in \emph{canonical ISV coverage}, $\Delta = \text{Primed}-\text{Direct}$.
Because each configuration cell has $n{=}3$ repeats, all quantitative reporting is descriptive: we do not perform significance testing and we do not claim causal identification of priming as a mechanism.

The study is deliberately comparative across regimes rather than optimized for a single leaderboard score.
HIGH asks whether priming is associated with coverage shifts when overlap is intentionally high.
LOW asks whether positive shifts can still appear when overlap is intentionally reduced.
UNSEEN asks whether positive shifts appear for a short technical/expository domain that is not represented in the narrative priming materials.
Together, the three regimes allow a restrained statement about persistence and heterogeneity without pretending that overlap, length, and genre have been orthogonally crossed.

\paragraph{Contributions.}
Conservatively stated, this paper contributes:
\begin{enumerate}[leftmargin=*]
  \item a reproducible Direct/Primed protocol for Polish-to-ISV LLM translation across seven configurations and three source regimes;
  \item a descriptive cross-regime evidence record for canonical coverage $\Delta$, including matched-configuration summaries and one documented invalid pairing cell;
  \item a qualitative paired audit synthesizing corpus-opening reuse, lexical borrowing, and orthographic change without equating metric gains with translation quality.
\end{enumerate}
We do \emph{not} claim that corpus priming universally improves translation, that higher overlap causes larger effects, or that any configuration is globally best.


\paragraph{Scope of claims.}
Throughout the paper we use associative language (``associated with'', ``observed'', ``descriptive'') rather than causal language (``improves'', ``causes'', ``proves''), except where we explicitly reject causal slogans.
This is not hedging for its own sake; it is required by the $n{=}3$ descriptive design and by the confounded regime factors.
